% GENERATED FILE. Edit canonical lessons, then run npm run generate:books.
% canonical-source-hash: fnv1a64:1e340ae1fc87165b
% canonical-lessons: PA-C351-gend, PA-C351-haki, PA-C351-tim, PA-C351-maich, PA-C351-harauna

\chapter{Ball Games}
\label{ch:pa-a2-ball-games}
\hlchaptermodality{351}

\begin{chapteropening}
\textbf{By the end of this chapter:} \emph{I can say which ball games I like to play.}

Complete the last lesson of chapter 351: I can say which ball games I like to play.
\end{chapteropening}

\section[\texorpdfstring{\pa{ਗੇਂਦ} (gend)}{ਗੇਂਦ (gend)}]{\pa{ਗੇਂਦ} (gend) — a ball}

\label{lesson:PA-C351-gend}



\begin{quote}
Before the new one: say the Punjabi for to make someone dance, then the Punjabi for to entertain.
\end{quote}

\subsection*{You'll want to know: \pa{ਗੇਂਦ}}
\textbf{\pa{ਗੇਂਦ}} — \emph{gend} — \textquotedblleft{}a ball\textquotedblright{}.

\begin{quote}
\pa{ਬੱਚੇ} \pa{ਬਾਹਰ} \pa{ਗੇਂਦ} \pa{ਨਾਲ} \pa{ਖੇਡਦੇ} \pa{ਹਨ।}
\end{quote}

\begin{quote}
\emph{bachche bāhar gend nāl kheḍde han.}
\end{quote}

\begin{quote}
The children play outside with a ball.
\end{quote}

A feminine noun.

\begin{cousinweb}
\textbf{\pa{ਗੇਂਦ}} (\emph{gend}) is inherited from Sanskrit \emph{genduka}, \textquotedblleft{}a ball for play\textquotedblright{}; CDIAL lists Punjabi \emph{gẽd}, \textquotedblleft{}ball\textquotedblright{}, under it (CDIAL 4248).
\end{cousinweb}

\begin{grammarlens}[title={What you've built}]
One more word for this chapter.
\end{grammarlens}

\subsection*{Your turn}
\begin{itemize}
\raggedright
  \item \emph{Say it:} \emph{gend}
  \item \emph{Say it:} \emph{gend}, once more
  \item \emph{Say it:} \emph{bachche bāhar gend nāl kheḍde han.} — The children play outside with a ball.
  \item \emph{Recall:} say the Punjabi for to make someone dance, then the Punjabi for to entertain, then say \emph{gend} again
\end{itemize}

\begin{tcolorbox}[breakable,colback=teal!4,colframe=teal!35!black,title={Before you move on}]
What does \pa{ਗੇਂਦ} mean? (\textquotedblleft{}A ball\textquotedblright{}.) Say it once more.
\end{tcolorbox}
\section[\texorpdfstring{\pa{ਹਾਕੀ} (hākī)}{ਹਾਕੀ (hākī)}]{\pa{ਹਾਕੀ} (hākī) — hockey}

\label{lesson:PA-C351-haki}



\begin{quote}
Before the new one: say the Punjabi for to entertain, then the Punjabi for a ball.
\end{quote}

\subsection*{You'll want to know: \pa{ਹਾਕੀ}}
\textbf{\pa{ਹਾਕੀ}} — \emph{hākī} — \textquotedblleft{}hockey\textquotedblright{}.

\begin{quote}
\pa{ਮੈਨੂੰ} \pa{ਹਾਕੀ} \pa{ਬਹੁਤ} \pa{ਪਸੰਦ} \pa{ਹੈ।}
\end{quote}

\begin{quote}
\emph{mainū̃ hākī bahut pasand hai.}
\end{quote}

\begin{quote}
I like hockey very much.
\end{quote}

A feminine noun.

\begin{cousinweb}
\textbf{\pa{ਹਾਕੀ}} (\emph{hākī}) is borrowed from English \emph{hockey}.
\end{cousinweb}

\begin{grammarlens}[title={What you've built}]
One more word for this chapter.
\end{grammarlens}

\subsection*{Your turn}
\begin{itemize}
\raggedright
  \item \emph{Say it:} \emph{hākī}
  \item \emph{Say it:} \emph{hākī}, once more
  \item \emph{Say it:} \emph{mainū̃ hākī bahut pasand hai.} — I like hockey very much.
  \item \emph{Recall:} say the Punjabi for to entertain, then the Punjabi for a ball, then say \emph{hākī} again
  \item \emph{Write it:} copy \textbf{\pa{ਹਾਕੀ}} once, checking the whole word
\end{itemize}

\begin{tcolorbox}[breakable,colback=teal!4,colframe=teal!35!black,title={Before you move on}]
What does \pa{ਹਾਕੀ} mean? (\textquotedblleft{}Hockey\textquotedblright{}.) Say it once more.
\end{tcolorbox}
\section[\texorpdfstring{\pa{ਟੀਮ} (ṭīm)}{ਟੀਮ (ṭīm)}]{\pa{ਟੀਮ} (ṭīm) — a team}

\label{lesson:PA-C351-tim}



\begin{quote}
Before the new one: say the Punjabi for a ball, then the Punjabi for hockey.
\end{quote}

\subsection*{You'll want to know: \pa{ਟੀਮ}}
\textbf{\pa{ਟੀਮ}} — \emph{ṭīm} — \textquotedblleft{}a team\textquotedblright{}.

\begin{quote}
\pa{ਸਾਡੀ} \pa{ਟੀਮ} \pa{ਬਹੁਤ} \pa{ਤਕੜੀ} \pa{ਹੈ।}
\end{quote}

\begin{quote}
\emph{sāḍī ṭīm bahut takṛī hai.}
\end{quote}

\begin{quote}
Our team is very strong.
\end{quote}

A feminine noun.

\begin{cousinweb}
\textbf{\pa{ਟੀਮ}} (\emph{ṭīm}) is borrowed from English \emph{team}.
\end{cousinweb}

\begin{grammarlens}[title={What you've built}]
One more word for this chapter.
\end{grammarlens}

\subsection*{Your turn}
\begin{itemize}
\raggedright
  \item \emph{Say it:} \emph{ṭīm}
  \item \emph{Say it:} \emph{ṭīm}, once more
  \item \emph{Say it:} \emph{sāḍī ṭīm bahut takṛī hai.} — Our team is very strong.
  \item \emph{Recall:} say the Punjabi for a ball, then the Punjabi for hockey, then say \emph{ṭīm} again
  \item \emph{Write it:} copy \textbf{\pa{ਟੀਮ}} once, checking the whole word
\end{itemize}

\begin{tcolorbox}[breakable,colback=teal!4,colframe=teal!35!black,title={Before you move on}]
What does \pa{ਟੀਮ} mean? (\textquotedblleft{}A team\textquotedblright{}.) Say it once more.
\end{tcolorbox}
\section[\texorpdfstring{\pa{ਮੈਚ} (maich)}{ਮੈਚ (maich)}]{\pa{ਮੈਚ} (maich) — a match (a game against a side)}

\label{lesson:PA-C351-maich}



\begin{quote}
Before the new one: say the Punjabi for hockey, then the Punjabi for a team.
\end{quote}

\subsection*{You'll want to know: \pa{ਮੈਚ}}
\textbf{\pa{ਮੈਚ}} — \emph{maich} — \textquotedblleft{}a match (a game against a side)\textquotedblright{}.

\begin{quote}
\pa{ਕੱਲ੍ਹ} \pa{ਸਾਡਾ} \pa{ਹਾਕੀ} \pa{ਦਾ} \pa{ਮੈਚ} \pa{ਹੋਵੇਗਾ।}
\end{quote}

\begin{quote}
\emph{kall sāḍā hākī dā maich hovegā.}
\end{quote}

\begin{quote}
Our hockey match will be tomorrow.
\end{quote}

A masculine noun.

\begin{cousinweb}
\textbf{\pa{ਮੈਚ}} (\emph{maich}) is borrowed from English \emph{match}.
\end{cousinweb}

\begin{grammarlens}[title={What you've built}]
One more word for this chapter.
\end{grammarlens}

\subsection*{Your turn}
\begin{itemize}
\raggedright
  \item \emph{Say it:} \emph{maich}
  \item \emph{Say it:} \emph{maich}, once more
  \item \emph{Say it:} \emph{kall sāḍā hākī dā maich hovegā.} — Our hockey match will be tomorrow.
  \item \emph{Recall:} say the Punjabi for hockey, then the Punjabi for a team, then say \emph{maich} again
  \item \emph{Write it:} copy \textbf{\pa{ਮੈਚ}} once, checking the whole word
\end{itemize}

\begin{tcolorbox}[breakable,colback=teal!4,colframe=teal!35!black,title={Before you move on}]
What does \pa{ਮੈਚ} mean? (\textquotedblleft{}A match (a game against a side)\textquotedblright{}.) Say it once more.
\end{tcolorbox}
\section[\texorpdfstring{\pa{ਹਰਾਉਣਾ} (harāuṇā)}{ਹਰਾਉਣਾ (harāuṇā)}]{\pa{ਹਰਾਉਣਾ} (harāuṇā) — to defeat; to beat (in a game)}

\label{lesson:PA-C351-harauna}



\begin{quote}
Before the new one: say the Punjabi for a team, then the Punjabi for a match.
\end{quote}

\subsection*{You'll want to know: \pa{ਹਰਾਉਣਾ}}
\textbf{\pa{ਹਰਾਉਣਾ}} — \emph{harāuṇā} — \textquotedblleft{}to defeat; to beat (in a game)\textquotedblright{}.

\begin{quote}
\pa{ਸਾਡੀ} \pa{ਟੀਮ} \pa{ਹਰ} \pa{ਮੈਚ} \pa{ਵਿੱਚ} \pa{ਉਸ} \pa{ਟੀਮ} \pa{ਨੂੰ} \pa{ਹਰਾਉਂਦੀ} \pa{ਹੈ।}
\end{quote}

\begin{quote}
\emph{sāḍī ṭīm har maich vich us ṭīm nū̃ harāundī hai.}
\end{quote}

\begin{quote}
Our team beats that team in every match.
\end{quote}

Past (I …): \pa{ਮੈਂ} \pa{ਅਮਨ} \pa{ਨੂੰ} \pa{ਹਰਾਇਆ}, \emph{maiṁ Aman nū̃ harāiā}.

\begin{cousinweb}
\textbf{\pa{ਹਰਾਉਣਾ}} (\emph{harāuṇā}) is inherited from Sanskrit \emph{hārayati}, \textquotedblleft{}causes to be taken, loses\textquotedblright{}; CDIAL lists Punjabi \emph{harāuṇā}, \textquotedblleft{}to defeat\textquotedblright{}, under it (CDIAL 14061).
\end{cousinweb}

\begin{grammarlens}[title={What you've built}]
One more word for this chapter.
\end{grammarlens}

\subsection*{Your turn}
\begin{itemize}
\raggedright
  \item \emph{Say it:} \emph{harāuṇā}
  \item \emph{Say it:} \emph{harāuṇā}, once more
  \item \emph{Say it:} \emph{sāḍī ṭīm har maich vich us ṭīm nū̃ harāundī hai.} — Our team beats that team in every match.
  \item \emph{Recall:} say the Punjabi for a team, then the Punjabi for a match, then say \emph{harāuṇā} again
  \item \emph{Write it:} copy \textbf{\pa{ਹਰਾਉਣਾ}} once, checking the whole word
\end{itemize}

\begin{tcolorbox}[breakable,colback=teal!4,colframe=teal!35!black,title={Before you move on}]
What does \pa{ਹਰਾਉਣਾ} mean? (\textquotedblleft{}To defeat; to beat (in a game)\textquotedblright{}.) Say it once more.
\end{tcolorbox}
